\documentclass[a4paper,12pt]{article}

\usepackage[english]{babel}
\usepackage[utf8]{inputenc}
\usepackage[T1]{fontenc}
\usepackage{geometry}
\geometry{margin=2.5cm}
\usepackage{hyperref}
\usepackage{microtype}
\usepackage{fancyhdr}
\usepackage{titlesec}
\usepackage{enumitem}
\usepackage{parskip}

\pagestyle{fancy}
\fancyhf{}
\lhead{Exam Answers}
\rhead{The Culture of Business}
\cfoot{\thepage}

\titleformat{\section}{\large\bfseries}{}{0em}{}[\titlerule]
\titleformat{\subsection}{\normalsize\bfseries}{}{0em}{}

\title{\vspace{-2cm}\textbf{Exam Answers}\\[0.2cm]
\large The Culture of Business, Finance and the Economy}
\author{Juan Ignacio Elosegui}
\date{Universidad Torcuato Di Tella \\ 2026}

\begin{document}
\maketitle

\noindent\textbf{Instructions:} Answer \textbf{4 questions} in total: at least 3 from questions 1--8 and no more than 1 from questions 9--11. This document prepares 5 answers: \textbf{Q2, Q3, Q4, Q5} (mandatory group) and \textbf{Q10} (elective group).

\vspace{1em}

% ─────────────────────────────────────────────
\section{Question 2 — The Faustian Bargain of Mathematics (Class IV)}

\textbf{Question:} \textit{Explain: in order to become a discipline that could ``leave no problem unsolved,'' as Viète promised, mathematics made a Faustian bargain.}

Mathematics began as \textbf{geometry}: divine, perfect, objective. Galileo's ``language of mathematics'' meant geometric diagrams. Geometry had strict limits --- irrationals like $\sqrt{2}$ were a conceptual scandal, and quantities of different kinds could not be mixed.

Alongside it existed a \textbf{proto-algebra}: dirty, plebeian, mercantile --- associated with traders and sailors, not universities. Algebra kept producing entities that had no geometric grounding: irrationals, and worse, imaginaries like $\sqrt{-1}$, which cannot even be approximated.

Viète and Oughtred made the decisive move: they replaced verbal descriptions with symbols, allowing algebra to operate on general classes of equations. Mathematics became conceivable as a universal language --- Viète's promise to ``leave no problem unsolved.''

The \textbf{Faustian bargain}: to gain that universality, mathematics had to (1) absorb the ``dirty'' mercantile algebra and (2) accept irrationals and imaginaries --- numbers that violated every classical ideal and corresponded to no object in the physical world. Mathematics sold its geometric soul to gain the world.

\vspace{1.5em}

% ─────────────────────────────────────────────
\section{Question 3 — Why and How Economics Adopted Mathematics (Class V)}

\textbf{Question:} \textit{Explain why, how, and what happened when Economics adopted mathematics.}

\textbf{Why:} Before the 1870s, the discipline was \textbf{political economy} --- closer to moral philosophy than science. Mill declared in 1848 that theory was ``complete.'' Reality proved otherwise: labour protests, rising public debt. The discipline needed to remake itself.

\textbf{How:} \textbf{Jevons} rebuilt economics on Bentham's utilitarianism (humans minimise pain, maximise pleasure) and argued it ``must be mathematical, simply because it deals with quantities.'' The tool: Leibniz's calculus --- the derivative --- which enabled the principle of \textbf{diminishing marginal utility}, the foundation of contemporary economics.

\textbf{What happened:} Everything that did not fit the calculus was expelled --- class struggle, moral economy, land ownership debates. Economics was freed from material knowledge: no need to know about farms, ships, or coal. Everything could be done from a desk.

\textbf{Why it stuck:} Marshall was cautious; Krugman said economists confused mathematical beauty with truth; Piketty called it ``infantile passion for mathematics.'' Yet mathematisation became entrenched: it functions as a rite of passage, confers status and technocratic authority (McCloskey: ``look how scientific I am''), and --- crucially --- economists \textit{created economic phenomena through mathematics}, making it impossible to know them any other way.

\vspace{1.5em}

% ─────────────────────────────────────────────
\section{Question 4 — Why and How Economics Modelled Itself on Physics (Class VI)}

\textbf{Question:} \textit{Explain why, how, and what happened when Economics modelled itself on physics.}

\textbf{Why:} This was a \textit{conscious} choice, not a natural convergence. Economists borrowed from physics to justify mathematical reasoning and claim scientific authority. Physics was the queen of the sciences --- light, wave theory, atoms, X-rays --- seemingly approaching a unified theory of everything. Gossen: ``Economics is concerned with the interaction of varying forces\ldots impossible to present without mathematics, just as with physics, mechanics.''

\textbf{How:} Three concrete cases:
\begin{itemize}
  \item \textbf{Fisher}: built a hydraulic mechanism using fluids as money to simulate equilibrium prices and quantities.
  \item \textbf{Phillips}: built the MONIAC, a hydraulic machine where coloured water represented income, spending, and GDP across tanks for households, firms, government, and trade.
  \item \textbf{Walras}: linked economics to Newton's laws, equating velocity/space/time with scarcity/utility/quantity.
\end{itemize}

Fisher's translation table persists to this day: particle = individual, space = commodity, energy = utility.

\textbf{What happened:} Jevons argued (1) utility, like gravity, is imperceptible but traceable through its effects (wages, prices); and (2) just as physics abstracts from complexity, economics need not know every market detail to believe in \textit{The Market}. The vocabulary reveals the lasting imprint: \textit{liquidity}, \textit{flows}, \textit{pipelines} --- words from physics that now sound entirely natural in economics.

\vspace{1.5em}

% ─────────────────────────────────────────────
\section{Question 5 — What Management Consultants ``Really'' Do (Class XI)}

\textbf{Question:} \textit{What do management consultants ``really'' do, how do they do it, and what skills matter most in their craft?}

Rahul, an Excel coach at a consulting boot camp: \textit{``All that's left after a client has paid millions of dollars is a presentation [of PowerPoint slides], about this thick.''} Consultants create, sustain, and reproduce \textbf{abstract work}.

Their work is defined by four separations: from ideas (they sell abstractions --- ``management,'' ``efficiency,'' ``impact''); from people (they change the work of people they never meet); from processes (they neither make nor sell the goods around which client firms are organised); and from the concrete (their work remains opaque even to themselves).

They are elite university graduates hired not for expertise but for ambition, hierarchy-respect, and capacity for long hours. Training is about speed: role-play exercises simulate a distracted CEO in an elevator, forcing findings into seconds.

\textbf{The paper mill case (Stein):} two weeks to find ``growth potential'' in a mill they know nothing about. Method: colour-code forklift drivers' activities as value-added (green), enabler (yellow), or non-value-added (red). Eliminate red = cut staff = ``improvement potential.'' Lucien, a French consultant writing a slide to sell insurance in Asia: \textit{``I know nothing about insurance; absolutely nothing''} --- yet convinced the client of his expertise.

\textbf{The real function:} consultants provide \textbf{legitimacy} both from what they \textit{don't} know (outsiders, free from internal loyalties) and what they \textit{do} know (they hold meetings with all managers, access confidential information, get employee badges). They are \textbf{decoupling technologies}: hired to push through decisions already made. A Munich interlocutor: \textit{``consulting firm XY said we should do this --- so it will be done.''} The key skill is not sector knowledge but producing abstractions that make the political appear technical.

\vspace{1.5em}

% ─────────────────────────────────────────────
\section{Question 10 — Excel and PowerPoint as Rhetorical Technologies (Class XII)}

\textbf{Question:} \textit{Discuss: ``Excel and PowerPoint are rhetorical technologies.``}

Final reports were slide decks. Consultants almost never used a word processor. PowerPoint was a major reason consulting work remained opaque: slides had to \textit{appear} intuitive, provoking automatic comprehension. Aesthetics mattered enormously --- hours spent aligning boxes, adjusting bullet styles. The only permitted complexity was mathematical; the mess behind the numbers stayed hidden.

Presentations followed strict narrative arcs: \textbf{situation, complication, solution}. The ``Mickey Mouse outline'' (mini-slides sketched like a comic strip) checked graphic and narrative flow before building the real deck. The naïveté of the visuals --- arrows hitting targets, a businessman scratching his head --- contrasted strikingly with the supposed weight of the work.

\textbf{Bullet points} are the basic cognitive unit: list-making, not argument. What matters is whether text groups under a heading, not logical relationships between points. Compared with prose, bullets make agreement easier --- \textit{``a language well suited to opacity that blurs disagreement''} --- freeing communication from logical coherence and allowing divergent understandings (what does ``resilience'' mean when my job is on the line?) to coexist without friction.

The \textbf{Change Matrix} (Chong) is the paradigmatic artefact: a visual displaying organisational states (``Burning platform,'' ``Building momentum'') that \textit{looks} precise but stays vague enough for divergent readings. Consultants impose shared vocabulary without forcing agreement on facts.

Excel and PowerPoint are rhetorical technologies because their primary function is not accurate description but \textbf{persuasion}: they produce an appearance of rigour that legitimises decisions already made, direct attention toward numbers and away from the mess, and make the political appear technical.

\end{document}
